\chapter{Model-View-*}

\section*{Tujuan Pembelajaran}

\begin{enumerate}
  \item Menjelaskan pembagian peran pada keempat varian Model-View-*, yaitu
    MVC, MVP, MVVM, dan MVI, beserta arah ketergantungan di antara komponennya.
  \item Membedakan keempat varian tersebut berdasarkan komponen yang dikenal
    oleh View, lalu menentukan varian mana yang sesuai untuk sebuah kebutuhan.
  \item Mengukur jumlah ketergantungan, kemudahan pengujian tanpa antarmuka,
    dan luas dampak perubahan pada keempat varian.
\end{enumerate}

\section{Pendahuluan}

Antarmuka pengguna berubah jauh lebih sering daripada aturan bisnis. Warna,
susunan tombol, dan urutan layar dapat berganti beberapa kali dalam satu
semester, sedangkan aturan perhitungannya tetap. Mencampur keduanya di dalam
satu berkas membuat setiap perubahan tampilan berisiko merusak perhitungan.

Kelompok pola Model-View-* lahir untuk memisahkan keduanya. Seluruh varian
membagi aplikasi menjadi bagian yang menyimpan data, bagian yang menampilkan,
dan bagian yang menghubungkan keduanya. Yang membedakan antar varian hanyalah
bagaimana bagian penghubung tersebut bekerja, dan komponen mana yang boleh
dikenal oleh View.

Bab ini membahas empat varian yang paling banyak dipakai, yaitu
Model-View-Controller (MVC), Model-View-Presenter (MVP),
Model-View-ViewModel (MVVM), dan Model-View-Intent (MVI). Keempatnya
diperbandingkan memakai satu aplikasi yang sama, sehingga perbedaannya dapat
diukur, bukan sekadar diperdebatkan. Fowler~\cite{fowler2006guiarch} menyusun
perbandingan sejenis dan memakai istilah \textit{passive view} untuk salah satu
bentuk yang dibahas di sini. Pembahasan mengikuti kerangka yang disusun
oleh Richards dan Ford~\cite{richards2020fundamentals}.

\section{Peran Ketiga Komponen}

Seluruh varian Model-View-* berangkat dari pembagian yang sama.
Tabel~\ref{tab:peran-mv} merangkum ketiga peran dasarnya.

\begin{table}[htbp]
  \centering
  \caption{Peran dasar yang dipakai seluruh varian Model-View-*}
  \label{tab:peran-mv}
  \begin{tabularx}{\textwidth}{|l|X|}
    \hline
    \textbf{Peran} & \textbf{Tanggung jawab} \\
    \hline
    Model & Menyimpan data dan menjalankan aturan bisnis \\
    \hline
    View & Menampilkan data dan menerima tindakan pengguna \\
    \hline
    Penghubung & Menerjemahkan tindakan pengguna menjadi perubahan pada Model \\
    \hline
  \end{tabularx}
\end{table}

Nama peran penghubung berbeda pada setiap varian. Pada MVC namanya Controller,
pada MVP namanya Presenter, pada MVVM namanya ViewModel, sedangkan pada MVI
peran tersebut dipegang oleh sebuah fungsi reduksi, bukan oleh sebuah kelas.
Perbedaan nama tersebut tidak penting. Yang menentukan adalah arah
ketergantungan di antara ketiganya.

\section{Empat Varian dan Arah Ketergantungannya}

Gambar~\ref{fig:empat-varian} menunjukkan keempat varian berdampingan. Panah
menyatakan arah ketergantungan, yaitu komponen yang menunjuk mengenal komponen
yang ditunjuk.

\begin{figure}[htbp]
  \centering
  \begin{tikzpicture}[
    kotak/.style={draw, rounded corners=2pt, minimum width=1.55cm,
      minimum height=0.62cm, align=center, font=\tiny\bfseries, inner sep=2pt},
    panah/.style={-{Stealth[length=1.6mm]}, thick},
    judul/.style={font=\scriptsize\bfseries}
  ]
    % MVC, pengguna menekan lewat Controller
    \node[kotak, fill=cView] (mvc_v) {View};
    \node[kotak, fill=cModel, below left=0.55cm and -0.5cm of mvc_v] (mvc_m) {Model};
    \node[kotak, fill=cPenghubung, above=0.55cm of mvc_v] (mvc_c) {Controller};
    \node[judul, above=0.2cm of mvc_c] {MVC};
    \draw[panah] (mvc_c) -- (mvc_v);
    \draw[panah] (mvc_c.west) -| (mvc_m.north);
    \draw[panah] (mvc_v) -- (mvc_m);

    % MVP
    \node[kotak, fill=cView, right=1.75cm of mvc_v] (mvp_v) {View};
    \node[kotak, fill=cModel, below left=0.55cm and -0.5cm of mvp_v] (mvp_m) {Model};
    \node[kotak, fill=cPenghubung, above=0.55cm of mvp_v] (mvp_p) {Presenter};
    \node[judul, above=0.2cm of mvp_p] {MVP};
    \draw[panah] (mvp_p) -- (mvp_v);
    \draw[panah] (mvp_p.west) -| (mvp_m.north);

    % MVVM
    \node[kotak, fill=cView, right=1.75cm of mvp_v, yshift=-1.17cm] (mvvm_v) {View};
    \node[kotak, fill=cPenghubung, above=0.55cm of mvvm_v] (mvvm_vm) {ViewModel};
    \node[kotak, fill=cModel, above=0.55cm of mvvm_vm] (mvvm_m) {Model};
    \node[judul, above=0.2cm of mvvm_m] {MVVM};
    \draw[panah] (mvvm_v) -- (mvvm_vm);
    \draw[panah] (mvvm_vm) -- (mvvm_m);

    % MVI
    \node[kotak, fill=cView, right=1.75cm of mvvm_v] (mvi_v) {View};
    \node[kotak, fill=cModel, above=0.55cm of mvi_v] (mvi_s) {Keadaan};
    \node[kotak, fill=cPenghubung, above=0.55cm of mvi_s] (mvi_r) {reduksi()};
    \node[judul, above=0.2cm of mvi_r] {MVI};
    \draw[panah] (mvi_r) -- (mvi_s);
    \draw[panah] (mvi_s) -- (mvi_v);
  \end{tikzpicture}
  \caption{Arah ketergantungan pada keempat varian. MVC di sini digambarkan
    dalam bentuk masa kini, sedangkan bentuk klasiknya berbeda dan dibahas pada
    Gambar~\ref{fig:mvc-klasik}. Setiap peran diberi warna
    tersendiri, yaitu ungu untuk Model, biru untuk penghubung, dan hijau untuk
    View. Letak pengguna dibahas pada gambar setiap varian}
  \label{fig:empat-varian}
\end{figure}

\subsection{Model-View-Controller}

MVC punya dua bentuk yang berbeda, dan keduanya sama-sama disebut MVC. Bentuk
pertama adalah bentuk asli dari Smalltalk-80~\cite{reenskaug1979mvc,krasner1988cookbook},
sedangkan bentuk kedua adalah yang
dipakai kerangka kerja web dan aplikasi desktop masa kini.
Gambar~\ref{fig:mvc-klasik} menunjukkan bentuk aslinya.

\begin{figure}[htbp]
  \centering
  \begin{tikzpicture}[
    kotak/.style={draw, rounded corners=2pt, minimum width=2.1cm,
      minimum height=0.7cm, align=center, font=\scriptsize\bfseries},
    panah/.style={-{Stealth[length=1.8mm]}, thick},
    lapor/.style={-{Stealth[length=1.8mm]}, thick, praditagreen, dashed}
  ]
    \node[kotak, fill=cPengguna] (u) {Pengguna};
    \node[kotak, fill=cPenghubung, right=1.5cm of u] (c) {Controller};
    \node[kotak, fill=cModel, right=1.8cm of c] (m) {Model};
    \node[kotak, fill=cView, below=1.3cm of m] (v) {View};

    \draw[panah] (u) -- node[above, font=\tiny] {menekan,}
      node[below, font=\tiny] {mengetik} (c);
    \draw[panah] (c) -- node[above, font=\tiny] {mengubah} (m);
    \draw[lapor] ([xshift=0.4cm]m.south) -- ([xshift=0.4cm]v.north)
      node[midway, right, font=\tiny] {memberi tahu};
    \draw[panah] ([xshift=-0.4cm]v.north) -- ([xshift=-0.4cm]m.south)
      node[midway, left, font=\tiny] {membaca};
    \draw[panah] (v.west) -- (u.south east)
      node[midway, below, font=\tiny] {menampilkan};
  \end{tikzpicture}
  \caption{MVC klasik. Controller mengubah Model, lalu Model memberi tahu View,
    dan View membaca Model untuk menggambar ulang}
  \label{fig:mvc-klasik}
\end{figure}

Pada bentuk klasik, alurnya berputar. Controller menerima tindakan pengguna lalu
mengubah Model. Model kemudian memberi tahu seluruh pihak yang mengamatinya,
termasuk View. View lalu membaca Model dan menggambar ulang dirinya. Model tidak
mengetahui jenis View yang mengamatinya, sebab pemberitahuannya memakai
mekanisme pengamatan.

Penelusuran satu tindakan memperjelas urutannya. Pengguna memasukkan 100 USD ke
IDR, lalu langkahnya sebagai berikut.

\begin{enumerate}
  \item View mendaftarkan diri sebagai pengamat Model. Pendaftaran tersebut
    dijalankan satu kali saat aplikasi disusun, bukan pada setiap tindakan.
  \item Pengguna menekan tombol, dan Controller yang menerimanya. Controller
    menafsirkan sendiri peristiwa tetikus dan papan ketik.
  \item Controller memanggil \texttt{model.konversi} dengan tiga nilai mentah.
  \item Model menghitung, menyimpan hasilnya sebagai \texttt{self.hasil}, lalu
    \textbf{memberi tahu seluruh pengamatnya}.
  \item View menerima pemberitahuan tersebut, lalu \textbf{menarik sendiri}
    nilai \texttt{model.hasil} dan menggambar ulang dirinya.
\end{enumerate}

Gambar~\ref{fig:sekuens-mvc-klasik} menunjukkan urutan tersebut sebagai diagram
sekuens.

\begin{figure}[htbp]
  \centering
  \includegraphics[width=0.95\textwidth]{sekuens-mvc-klasik.pdf}
  \caption{Diagram sekuens MVC klasik. Model yang memicu penggambaran, dan
    Controller tidak menyentuh View sama sekali}
  \label{fig:sekuens-mvc-klasik}
\end{figure}

Langkah keempat memuat perbedaan yang paling menentukan. Yang memicu
penggambaran adalah Model, bukan Controller. Controller bahkan tidak mengetahui
bahwa View sudah menggambar ulang. Susunan seperti itu memungkinkan satu Model
diamati banyak View sekaligus, dan seluruhnya diperbarui oleh satu pemberitahuan.

Bentuk kedua muncul karena kerangka kerja web bekerja dengan cara yang berbeda.
Permintaan datang, dilayani, lalu selesai, sehingga tidak ada View yang menetap
untuk diberi tahu. Gambar~\ref{fig:mvc-modern} menunjukkan bentuk tersebut.

Letak pengguna pada bentuk kedua berbeda dari bentuk klasiknya, dan perbedaan
itu sering terlewat. Pada bentuk klasik, Controller menafsirkan sendiri
peristiwa tetikus dan papan ketik, sehingga pengguna berhadapan langsung
dengannya. Pada bentuk masa kini, Controller tidak punya permukaan yang dapat
disentuh sama sekali.

Yang disentuh pengguna adalah View. Pada aplikasi web, pengguna menekan tautan
pada halaman yang sudah tergambar, lalu perute mengubahnya menjadi permintaan
dan mengantarkannya ke Controller. Pada aplikasi desktop, widget dimiliki View,
dan penangan widget itulah yang memanggil Controller. Keduanya berbeda pada
perantaranya, tetapi sama pada intinya, yaitu pengguna menyentuh View sedangkan
Controller hanya menerima permintaan yang sudah diteruskan.

\begin{figure}[htbp]
  \centering
  \begin{tikzpicture}[
    kotak/.style={draw, rounded corners=2pt, minimum width=2.1cm,
      minimum height=0.7cm, align=center, font=\scriptsize\bfseries},
    panah/.style={-{Stealth[length=1.8mm]}, thick}
  ]
    \node[kotak, fill=cPenghubung] (c) at (0,0) {Controller};
    \node[kotak, fill=cView] (v) at (0,-1.6) {View};
    \node[kotak, fill=cModel] (m) at (0,-3.2) {Model};
    \node[kotak, fill=cPengguna] (u) at (-3.7,-1.6) {Pengguna};
    % Pengguna dan View saling berhadapan.
    \draw[panah] ([yshift=1.6mm]u.east) -- ([yshift=1.6mm]v.west)
      node[midway, above, font=\tiny] {menekan};
    \draw[panah] ([yshift=-1.6mm]v.west) -- ([yshift=-1.6mm]u.east)
      node[midway, below, font=\tiny] {tampilan};
    % View meneruskan ke atas, Controller memerintah ke bawah.
    \draw[panah] ([xshift=-5mm]v.north) -- ([xshift=-5mm]c.south)
      node[midway, left, font=\tiny] {meneruskan};
    \draw[panah] ([xshift=5mm]c.south) -- ([xshift=5mm]v.north)
      node[midway, right, font=\tiny] {memerintah};
    % Controller menyerahkan rujukan Model, lalu View membaca hasilnya.
    \draw[panah] ([xshift=-5mm]v.south) -- ([xshift=-5mm]m.north)
      node[midway, left, font=\tiny] {membaca hasil};
    \draw[panah] (c.east) -- ++(1.7,0) |- (m.east)
      node[pos=0.22, right, font=\tiny] {memanggil};
  \end{tikzpicture}
  \caption{MVC pada kerangka kerja masa kini. Pengguna menyentuh View, lalu
    permintaannya diteruskan ke Controller. Controller memanggil Model, lalu
    memerintah View. View membaca hasil dari Model yang rujukannya sudah
    diserahkan Controller, dan Model tidak memberi tahu siapa pun}
  \label{fig:mvc-modern}
\end{figure}

Tabel~\ref{tab:dua-mvc} merangkum perbedaan keduanya.

\begin{table}[htbp]
  \centering
  \caption{Perbedaan MVC klasik dan MVC pada kerangka kerja masa kini}
  \label{tab:dua-mvc}
  \begin{tabularx}{\textwidth}{|l|X|X|}
    \hline
    \textbf{Hal} & \textbf{MVC klasik} & \textbf{MVC masa kini} \\
    \hline
    Pemicu penggambaran & Model memberi tahu View & Controller memerintah View \\
    \hline
    Sifat & Menetap, View hidup terus & Sekali jalan, permintaan lalu selesai \\
    \hline
    Peran Controller & Menangani masukan saja & Menangani masukan dan mengatur alur \\
    \hline
    Yang disentuh pengguna & Controller, sebab Controller menafsirkan tetikus dan papan ketik & View, lalu diteruskan ke Controller \\
    \hline
    Pengetahuan Model & Tahu ada yang mengamati & Tidak tahu apa pun \\
    \hline
  \end{tabularx}
\end{table}

Contoh pada bab ini memakai bentuk masa kini, sebab bentuk itulah yang paling
sering ditemui. Ciri khasnya terletak pada konstruktor View, yang menerima Model
sebagai argumen. Karena itu View tidak dapat dibuat tanpa Model, dan pengujian
logika tampilan tetap menuntut kedua objek tersebut ada. Tombolnya pun dimiliki
View, sehingga pengguna menyentuh View lalu View meneruskannya ke Controller.

\begin{lstlisting}[style=PythonStyle, caption={View pada MVC mengenal Model,
  sekaligus meneruskan tindakan pengguna, tersedia di
  \berkas{sources/chapter03/mvc.py}}, label=lst:mvc]
class View:

  def __init__(self, induk, model):
    self.model = model
    self.controller = None
    self.tombol = tk.Button(induk, text="Hitung",
                            command=self.saat_ditekan)

  def saat_ditekan(self):
    # Widget dimiliki View, sehingga pengguna menyentuh View.
    self.controller.tangani_masukan("USD", "IDR", 100)
\end{lstlisting}

Penelusuran satu permintaan memperjelas apa yang menyeberang antar komponen.
Pengguna memasukkan 100 USD ke IDR, lalu urutannya sebagai berikut.

\begin{enumerate}
  \item Pengguna menekan tombol milik View, lalu View memanggil
    \texttt{controller.tangani\_masukan} dengan tiga nilai mentah, yaitu
    \texttt{"USD"}, \texttt{"IDR"}, dan \texttt{100}.
  \item Controller memanggil \texttt{model.konversi} dengan ketiga nilai
    tersebut.
  \item Model menghitung, lalu \textbf{menyimpan hasilnya di dalam dirinya
    sendiri} sebagai \texttt{self.hasil}.
  \item Controller memanggil \texttt{view.render} tanpa membawa data apa pun.
  \item View \textbf{menarik sendiri} nilai \texttt{model.hasil}, lalu menyusun
    teks \texttt{Hasil: 1,625,000.00}.
\end{enumerate}

Gambar~\ref{fig:sekuens-mvc} menunjukkan urutan tersebut sebagai diagram
sekuens.

\begin{figure}[htbp]
  \centering
  \includegraphics[width=0.95\textwidth]{sekuens-mvc.pdf}
  \caption{Diagram sekuens MVC. Hasil disimpan di Model, lalu ditarik sendiri oleh View}
  \label{fig:sekuens-mvc}
\end{figure}

Ciri khas MVC terletak pada langkah kelima. Hasil tidak dikirimkan kepada View,
melainkan diambil View dari Model. Karena itu View wajib mengenal Model.

Perlu satu penegasan agar tidak salah paham. View memang membaca Model, tetapi
View tidak mencari sendiri Model yang mana. Rujukannya diserahkan lebih dahulu,
yaitu lewat konstruktor pada aplikasi desktop atau lewat objek yang dititipkan
Controller pada aplikasi web. Jadi Controller tetap memegang kendali alurnya,
sedangkan View hanya membaca isi objek yang sudah diberikan kepadanya.

\subsection{Model-View-Presenter}

MVP diperkenalkan Potel~\cite{potel1996mvp} sebagai penyederhanaan MVC. MVP
menjadikan View pasif. Presenter mengerjakan seluruh penyusunan teks, lalu
mendorong hasilnya ke View. Gambar~\ref{fig:mvp} menunjukkan bahwa tidak ada
panah dari View ke Model. Ruang di bawah View kosong, dan kekosongan itulah
pembeda MVP dari MVC.

\begin{figure}[htbp]
  \centering
  \begin{tikzpicture}[
    kotak/.style={draw, rounded corners=2pt, minimum width=2.1cm,
      minimum height=0.7cm, align=center, font=\scriptsize\bfseries},
    panah/.style={-{Stealth[length=1.8mm]}, thick}
  ]
    \node[kotak, fill=cPenghubung] (p) at (0,0) {Presenter};
    \node[kotak, fill=cView] (v) at (0,-1.6) {View};
    \node[kotak, fill=cModel] (m) at (0,-3.2) {Model};
    \node[kotak, fill=cPengguna] (u) at (-3.7,-1.6) {Pengguna};
    % Pengguna hanya berhadapan dengan View, tidak dengan Presenter.
    \draw[panah] ([yshift=1.6mm]u.east) -- ([yshift=1.6mm]v.west)
      node[midway, above, font=\tiny] {menyentuh};
    \draw[panah] ([yshift=-1.6mm]v.west) -- ([yshift=-1.6mm]u.east)
      node[midway, below, font=\tiny] {tampilan};
    % View meneruskan ke atas, Presenter mendorong teks ke bawah.
    \draw[panah] ([xshift=-5mm]v.north) -- ([xshift=-5mm]p.south)
      node[midway, left, font=\tiny] {meneruskan};
    \draw[panah] ([xshift=5mm]p.south) -- ([xshift=5mm]v.north)
      node[midway, right, font=\tiny] {mendorong teks};
    % Tidak ada panah dari View ke Model. Itulah inti passive view.
    \draw[panah] (p.east) -- ++(1.7,0) |- (m.east)
      node[pos=0.22, right, font=\tiny] {memanggil};
  \end{tikzpicture}
  \caption{MVP, dengan View pasif yang tidak mengenal Model. Pengguna tetap
    menyentuh View, lalu View meneruskannya ke Presenter}
  \label{fig:mvp}
\end{figure}

View pada MVP hanya menyimpan teks yang diterimanya. Karena tidak ada perilaku
yang perlu ditiru selain menyimpan, View tersebut dapat diganti tiruan sederhana
saat diuji.

Bentuk yang dipakai di sini disebut \textit{passive view}, yaitu View yang sama
sekali tidak mengenal komponen lain. Banyak penerapan MVP memakai bentuk lain,
yaitu View menyimpan rujukan ke Presenter agar dapat meneruskan tindakan
pengguna. Pada bentuk tersebut hubungannya menjadi dua arah, dan angka
ketergantungannya bertambah satu.

\begin{lstlisting}[style=PythonStyle, caption={View pasif pada MVP, tersedia di
  \berkas{sources/chapter03/mvp.py}}, label=lst:mvp]
class View:
  """Menyimpan teks yang diberikan Presenter, tanpa logika apa pun."""

  def tampilkan(self, teks):
    self.teks = teks
\end{lstlisting}

Penelusuran permintaan yang sama pada MVP berbeda mulai langkah ketiga.

\begin{enumerate}
  \item Presenter menerima tiga nilai mentah.
  \item Presenter memanggil \texttt{model.konversi}, dan Model
    \textbf{mengembalikan} angka \texttt{1625000.0} sebagai nilai balik.
  \item Presenter menyusun teks \texttt{Hasil: 1,625,000.00}.
  \item Presenter memanggil \texttt{view.tampilkan} dengan \textbf{membawa teks
    yang sudah jadi}.
  \item View menyimpan teks tersebut tanpa mengolahnya.
\end{enumerate}

Gambar~\ref{fig:sekuens-mvp} menunjukkan urutan tersebut sebagai diagram
sekuens.

\begin{figure}[htbp]
  \centering
  \includegraphics[width=0.95\textwidth]{sekuens-mvp.pdf}
  \caption{Diagram sekuens MVP. Teks yang sudah jadi didorong Presenter ke View}
  \label{fig:sekuens-mvp}
\end{figure}

Perbedaannya jelas. Pada MVC hasil disimpan di Model lalu ditarik View,
sedangkan pada MVP hasil dikembalikan ke Presenter lalu didorong ke View. View
karena itu tidak perlu mengenal Model.

\subsection{Model-View-ViewModel}

MVVM diperkenalkan Gossman~\cite{gossman2005mvvm} bersama kerangka kerja
Windows Presentation Foundation. MVVM membalik arah pemberitahuan. ViewModel
menyediakan properti yang dapat
diamati, lalu View mendaftarkan diri sebagai pengamat. Ketika properti berubah,
View memperbarui dirinya sendiri. Gambar~\ref{fig:mvvm} menunjukkan bahwa
ViewModel tidak mengenal View.

\begin{figure}[htbp]
  \centering
  \begin{tikzpicture}[
    kotak/.style={draw, rounded corners=2pt, minimum width=2.1cm,
      minimum height=0.7cm, align=center, font=\scriptsize\bfseries},
    panah/.style={-{Stealth[length=1.8mm]}, thick}
  ]
    \node[kotak, fill=cModel] (m) {Model};
    \node[kotak, fill=cPenghubung, below=0.6cm of m] (vm) {ViewModel};
    \node[kotak, fill=cView, below=0.6cm of vm] (v) {View};
    \node[kotak, fill=cPengguna, left=1.5cm of v] (u) {Pengguna};
    \draw[panah] (vm) -- node[right, font=\tiny] {memanggil} (m);
    \draw[panah] (v) -- node[right, font=\tiny] {mengamati} (vm);
    \draw[panah] (u) -- node[above, font=\tiny] {menyentuh} (v);
  \end{tikzpicture}
  \caption{MVVM, dengan View yang mengamati properti milik ViewModel}
  \label{fig:mvvm}
\end{figure}

Jumlah ketergantungan MVVM sama dengan MVC, yaitu tiga. Arahnya yang berbeda.
Pada MVC View mengenal Model, sedangkan pada MVVM View mengenal ViewModel dan
ViewModel tidak mengenal View. Karena itu logika tampilan dapat diuji cukup
dengan membaca properti, tanpa membuat objek View.

\begin{lstlisting}[style=PythonStyle, caption={View mengikat diri pada properti,
  tersedia di \berkas{sources/chapter03/mvvm.py}}, label=lst:mvvm]
class View:
  """Mengikat diri pada properti ViewModel dan menyimpan nilai terbarunya."""

  def __init__(self, view_model):
    self.terakhir = None
    view_model.teks.amati(self.saat_berubah)
\end{lstlisting}

Penelusuran permintaan yang sama pada MVVM berbeda pada arah pemberitahuannya.

\begin{enumerate}
  \item ViewModel menerima tiga nilai mentah.
  \item ViewModel memanggil \texttt{model.konversi} dan menerima
    \texttt{1625000.0}.
  \item ViewModel menulis teks ke propertinya lewat \texttt{teks.ubah}.
  \item Properti tersebut \textbf{memanggil seluruh pengamatnya} satu per satu.
  \item View menerima teks baru lewat \texttt{saat\_berubah}, lalu memperbarui
    dirinya.
\end{enumerate}

Gambar~\ref{fig:sekuens-mvvm} menunjukkan urutan tersebut sebagai diagram
sekuens. Perhatikan langkah pertama, yaitu View mendaftarkan diri sebagai
pengamat sebelum permintaan datang.

\begin{figure}[htbp]
  \centering
  \includegraphics[width=0.95\textwidth]{sekuens-mvvm.pdf}
  \caption{Diagram sekuens MVVM. Properti yang memberi tahu pengamatnya, bukan ViewModel}
  \label{fig:sekuens-mvvm}
\end{figure}

ViewModel tidak pernah memanggil View. ViewModel hanya mengubah nilai properti,
dan properti itulah yang memberi tahu siapa pun yang mendaftar. Karena itu satu
properti dapat diamati beberapa View sekaligus tanpa ViewModel mengetahuinya.

\subsection{Model-View-Intent}

MVI menghapus kelas penghubung sama sekali. Setiap tindakan pengguna dibungkus
menjadi Intent, lalu sebuah fungsi reduksi mengubah pasangan keadaan lama dan
Intent menjadi keadaan baru yang tidak dapat diubah lagi.
Gambar~\ref{fig:mvi} menunjukkan susunannya.

\begin{figure}[htbp]
  \centering
  \begin{tikzpicture}[
    kotak/.style={draw, rounded corners=2pt, minimum width=1.9cm,
      minimum height=0.7cm, align=center, font=\scriptsize\bfseries},
    panah/.style={-{Stealth[length=1.8mm]}, thick}
  ]
    \node[kotak, fill=cPenghubung] (i) {Intent};
    \node[kotak, fill=cPenghubung, right=1.0cm of i] (r) {reduksi()};
    \node[kotak, fill=cModel, right=1.0cm of r] (k) {Keadaan};
    \node[kotak, fill=cView, right=1.0cm of k] (v) {View};
    \node[kotak, fill=cPengguna, above=0.5cm of v] (u) {Pengguna};
    \draw[panah] (i) -- (r);
    \draw[panah] (r) -- (k);
    \draw[panah] (k) -- (v);
    \draw[panah] (u) -- node[right, font=\tiny] {menyentuh} (v);
    \draw[panah] (v.south) -- ++(0,-0.8) -| (i.south)
      node[pos=0.25, below, font=\tiny] {dibungkus menjadi Intent};
  \end{tikzpicture}
  \caption{MVI bersifat searah dan membentuk siklus. View menggambar Keadaan,
    lalu tindakan pengguna kembali menjadi Intent}
  \label{fig:mvi}
\end{figure}

Listing~\ref{lst:reduksi} memperlihatkan fungsi tersebut. Karena fungsi itu
murni, seluruh keadaan yang pernah digambar dapat disimpan lalu diputar ulang
persis seperti semula. Sifat itu membuat antarmuka yang keluarannya
berubah-ubah tetap dapat diuji secara pasti.

\begin{lstlisting}[style=PythonStyle, caption={Fungsi reduksi pada MVI,
  tersedia di \berkas{sources/chapter03/mvi.py}}, label=lst:reduksi]
def reduksi(keadaan, intent):
  """Mengubah keadaan lama dan sebuah Intent menjadi keadaan baru.

  Fungsi ini murni, artinya keluarannya hanya bergantung pada kedua masukannya.
  Sifat itulah yang membuat riwayat tampilan dapat diputar ulang.
  """
  try:
    hasil = domain.hitung_konversi(intent.kode_asal, intent.kode_tujuan,
                                   intent.nominal)
    return replace(keadaan, teks=f"Hasil: {hasil:,.2f}", sedang_sibuk=False)
  except (ValueError, LookupError) as kesalahan:
    return replace(keadaan, teks=f"Galat: {kesalahan}", sedang_sibuk=False)
\end{lstlisting}

Agar keempat varian tidak berhenti sebagai nama, Tabel~\ref{tab:contoh-mv}
menunjukkan isi setiap peran pada aplikasi konversi yang dipakai sepanjang bab
ini. Perhatikan bahwa isi Model dan View hampir sama pada keempatnya, sedangkan
peran penghubungnya berbeda jauh.

\begin{table}[htbp]
  \centering
  \caption{Contoh isi setiap peran pada aplikasi konversi}
  \label{tab:contoh-mv}
  \begin{tabularx}{\textwidth}{|l|X|X|}
    \hline
    \textbf{Varian} & \textbf{Isi penghubungnya} & \textbf{Yang dilakukan View} \\
    \hline
    MVC & Controller meneruskan masukan ke Model & Membaca Model lalu menyusun teks \\
    \hline
    MVP & Presenter menyusun teks lalu mendorongnya & Menyimpan teks yang diterima \\
    \hline
    MVVM & ViewModel menulis hasil ke properti & Memperbarui diri saat properti berubah \\
    \hline
    MVI & Fungsi reduksi menghasilkan keadaan baru & Menggambar keadaan yang diterima \\
    \hline
  \end{tabularx}
\end{table}

Kasus berikut memperlihatkan akibatnya pada pekerjaan sehari-hari. Sebuah tim
ingin menguji aturan penolakan nominal nol. Pada MVC, pengujian tersebut tetap
menuntut objek View dibuat lebih dahulu, sebab Controller merujuk View. View
menuntut Model, dan Model menuntut sambungan data. Pengujian yang seharusnya
memeriksa satu aturan akhirnya menyalakan hampir seluruh aplikasi. Pada MVI,
pengujian yang sama cukup memanggil satu fungsi reduksi tanpa membuat objek apa
pun.

Yang menentukan bukan jumlah kelasnya, melainkan komponen apa yang dituntut
View. Itulah sebabnya perbandingan pada bab ini memakai ketergantungan sebagai
ukuran, bukan panjang kode.

Penelusuran permintaan yang sama pada MVI tidak memakai pemanggilan metode
antar objek sama sekali.

\begin{enumerate}
  \item Tindakan pengguna dibungkus menjadi
    \texttt{IntentKonversi("USD", "IDR", 100)}.
  \item Fungsi \texttt{reduksi} dipanggil dengan dua masukan, yaitu keadaan lama
    dan Intent tersebut.
  \item Di dalamnya \texttt{domain.hitung\_konversi} menghasilkan
    \texttt{1625000.0}.
  \item Fungsi tersebut \textbf{mengembalikan objek Keadaan baru} berisi teks
    hasil, tanpa mengubah keadaan lama.
  \item View menggambar Keadaan tersebut, lalu menyimpannya ke dalam riwayat.
\end{enumerate}

Gambar~\ref{fig:sekuens-mvi} menunjukkan urutan tersebut sebagai diagram
sekuens.

\begin{figure}[htbp]
  \centering
  \includegraphics[width=0.95\textwidth]{sekuens-mvi.pdf}
  \caption{Diagram sekuens MVI. Setiap langkah menghasilkan objek baru, tanpa mengubah yang lama}
  \label{fig:sekuens-mvi}
\end{figure}

Seluruh isi tampilan terbungkus dalam satu objek Keadaan. Karena keadaan lama
tidak pernah diubah, seluruh riwayat tersimpan utuh dan dapat diputar ulang.

Tabel~\ref{tab:aliran-data} merangkum apa yang menyeberang menuju View pada
keempat varian.

\begin{table}[htbp]
  \centering
  \caption{Apa yang menyeberang menuju View pada keempat varian}
  \label{tab:aliran-data}
  \begin{tabularx}{\textwidth}{|l|X|X|}
    \hline
    \textbf{Varian} & \textbf{Yang diterima View} & \textbf{Cara View memperolehnya} \\
    \hline
    MVC & Tidak ada, hanya perintah menggambar & Menarik sendiri dari Model \\
    \hline
    MVP & Teks yang sudah jadi & Didorong Presenter lewat pemanggilan metode \\
    \hline
    MVVM & Nilai properti yang baru & Diberi tahu properti yang diamatinya \\
    \hline
    MVI & Satu objek Keadaan yang utuh & Diserahkan sebagai argumen saat digambar \\
    \hline
  \end{tabularx}
\end{table}

\section{Produk yang Memakai Setiap Varian}

Keempat varian tersebut bukan bahan kuliah semata. Setiap varian punya kerangka
kerja atau produk yang memakainya. Tabel~\ref{tab:produk-mv} merangkum contoh
yang paling sering ditemui.

\begin{table}[htbp]
  \centering
  \caption{Contoh produk dan kerangka kerja untuk setiap varian}
  \label{tab:produk-mv}
  \begin{tabularx}{\textwidth}{|l|X|X|}
    \hline
    \textbf{Varian} & \textbf{Contoh produk} & \textbf{Catatan} \\
    \hline
    MVC klasik & Smalltalk-80 & Bentuk aslinya, dengan Model yang memberi tahu View \\
    \hline
    MVC masa kini & Ruby on Rails, Laravel, Spring MVC, ASP.NET Core MVC & Controller menerima permintaan dari perute, lalu memerintah View \\
    \hline
    MVP & Google Web Toolkit, Windows Forms, aplikasi Android sebelum era Jetpack & Presenter mendorong teks jadi ke View yang pasif \\
    \hline
    MVVM & Windows Presentation Foundation, .NET MAUI, Vue.js versi awal & View mengamati properti pada ViewModel \\
    \hline
    MVI & Redux, arsitektur Elm, Flutter Bloc & Keadaan diganti utuh oleh fungsi reduksi. Android menganjurkan aliran searah dengan nama UDF, tanpa memakai istilah MVI \\
    \hline
  \end{tabularx}
\end{table}

Empat hal perlu penjelasan tambahan, sebab keempatnya sering menimbulkan salah paham.

\textbf{Django tidak memakai penamaan MVC.} Django menyebut polanya
\textit{Model-Template-View} atau MTV. Berkas yang disebut \texttt{views.py} pada
Django sebenarnya berperan sebagai Controller, sebab berkas tersebut memilih
data mana yang disajikan. Templat Django itulah yang berperan sebagai View.
Dokumentasi resminya menyatakan Controller pada Django adalah kerangka kerjanya
sendiri, yaitu bagian yang mengantar permintaan ke fungsi yang
tepat~\cite{django2026faq}.

\textbf{Android tidak menyebut dirinya MVVM.} Banyak tulisan menyebut Android
memakai MVVM karena ada kelas bernama \texttt{ViewModel}. Dokumentasi resminya
tidak memakai istilah tersebut sama sekali. Dokumentasi tersebut menyebut
\texttt{ViewModel} sebagai \textit{state holder}, yaitu pemegang keadaan layar,
di dalam susunan berlapis yang terdiri atas lapisan antarmuka, domain, dan
data~\cite{android2026viewmodel}. Kemiripannya dengan MVVM memang ada, sebab
View mengamati keadaan dan \texttt{ViewModel} tidak mengenal View. Meskipun
begitu, penyebutannya sebaiknya tetap hati-hati.

\textbf{Intent pada Android dan Intent pada MVI memakai gagasan yang sama.}
Kesamaan namanya bukan kebetulan. Keduanya membungkus niat menjadi objek data,
lalu menyerahkan objek tersebut kepada pihak yang menyalurkannya.

Intent untuk mengirim pesan singkat memperjelas hal tersebut. Objek tersebut
tidak mengirim apa pun sendiri. Objek tersebut hanya menyatakan adanya niat
mengirim pesan ke nomor tertentu. Sistem operasi lalu memilih aplikasi mana yang
mengerjakannya, dan pengguna bahkan dapat ikut memilih ketika aplikasinya lebih
dari satu. Dokumentasi Android menyebutnya objek pesan untuk meminta tindakan
dari komponen aplikasi lain~\cite{android2026intent}. Intent pada MVI bekerja
dengan cara yang sama, hanya saja penyalurnya berupa fungsi reduksi.

Tujuan keduanya pun sama. Sebuah Intent dikirim agar ada keadaan yang berubah,
atau agar ada proses tertentu yang dijalankan. Objek itu sendiri tidak
mengerjakan apa pun.

Gagasan tersebut dikenal luas dengan nama \textit{command}. Sebuah
\textit{command} adalah permintaan agar sistem melakukan tindakan yang mengubah
keadaannya. Wujudnya berupa kelas yang hanya berisi data yang diperlukan,
semacam \textit{Data Transfer Object} khusus~\cite{microsoft2021command}.
Rumusan tersebut berlaku untuk keduanya tanpa perubahan sedikit pun. Akibat yang
menyertainya juga sama. Pengirim tidak perlu mengenal pihak yang mengerjakan,
permintaannya dapat dicatat dan diperiksa, dan pengerjanya dapat diganti tanpa
menyentuh pengirimnya.

Yang berbeda hanyalah tingkat penerapannya, bukan gagasan maupun tujuannya.
Dua baris pertama pada Tabel~\ref{tab:dua-intent} karena itu berisi hal yang
persis sama untuk kedua kolomnya.

\begin{table}[htbp]
  \centering
  \caption{Intent pada Android dan Intent pada MVI, satu gagasan pada dua
    tingkat penerapan}
  \label{tab:dua-intent}
  \begin{tabularx}{\textwidth}{|l|X|X|}
    \hline
    \textbf{Hal} & \textbf{Intent pada Android} & \textbf{Intent pada MVI} \\
    \hline
    Gagasannya & Niat dibungkus menjadi objek data & Niat dibungkus menjadi objek data \\
    \hline
    Tujuannya & Ada keadaan yang berubah atau proses yang dijalankan & Ada keadaan yang berubah atau proses yang dijalankan \\
    \hline
    Contohnya & Niat mengirim pesan singkat ke satu nomor & Niat mengonversi 100 USD menjadi IDR \\
    \hline
    Penyalurnya & Sistem operasi Android & Fungsi reduksi \\
    \hline
    Cakupannya & Melintasi komponen dan aplikasi & Berada di dalam satu layar \\
    \hline
    Wujud perubahannya & Komponen lain dijalankan, misalnya aplikasi pesan terbuka & Objek keadaan baru dikembalikan \\
    \hline
  \end{tabularx}
\end{table}

Kesimpulannya perlu ditarik hati-hati. Pemakaian kelas \texttt{Intent} tidak
dengan sendirinya menjadikan sebuah aplikasi Android memakai MVI. MVI menuntut
lebih dari sekadar membungkus niat, yaitu keadaan layar yang diganti utuh oleh
fungsi reduksi. Bagian Android yang benar-benar sejajar dengan MVI justru
terletak di tempat lain. Dokumentasi Android menganjurkan
\textit{unidirectional data flow} atau UDF, yaitu aliran data satu arah.
Keadaan mengalir turun menuju antarmuka, sedangkan peristiwa mengalir naik
menuju \texttt{ViewModel}~\cite{android2026uilayer}. Bentuk aliran tersebut sama
dengan MVI, meskipun istilah MVI tidak dipakai dalam dokumentasinya.

\textbf{Redux mewarisi arsitektur Elm.} Dokumentasi Redux menyebut Elm sebagai
sumber gagasannya. Elm memakai susunan \textit{model view update}, dan fungsi
pembaruannya bertanda tangan \texttt{(action, state) => state}. Fungsi tersebut
sama perannya dengan fungsi reduksi pada Redux~\cite{redux2026priorart}. Karena
itu Redux menjadi contoh MVI yang paling mudah ditemui.

Vue.js pada dokumentasi versi awalnya menyebut setiap instansnya sebagai
ViewModel, dan karena itu variabelnya dinamai \texttt{vm}~\cite{vue2016instance}.
Dokumentasi versi terbarunya tidak lagi memakai istilah tersebut. Penyebutan
seperti ini memang berubah seiring waktu, sehingga rujukan versi perlu ikut
dicatat.

\section{Ketergantungan View sebagai Pembeda}

Satu pertanyaan cukup untuk membedakan keempat varian, yaitu komponen apa yang
dikenal oleh View. Jawabannya menentukan seberapa mudah logika tampilan diuji
tanpa menjalankan antarmuka sungguhan.

Analogi yang mendekati adalah loket informasi di kampus. Petugas loket yang
harus membuka lemari arsip sendiri sulit digantikan, sebab penggantinya harus
mengenal isi lemari tersebut. Petugas yang hanya membacakan lembar yang sudah
disiapkan bagian lain mudah digantikan siapa pun. View pada MVP dan MVI berperan
seperti petugas kedua.

Tabel~\ref{tab:ketergantungan-mv} merangkum hasil pemeriksaan nyata memakai
\texttt{periksa\_mv.py}. Kolom terakhir menunjukkan cara pengujian yang paling
sedikit menyentuh antarmuka.

\begin{table}[htbp]
  \centering
  \caption{Ketergantungan dan cara pengujian keempat varian}
  \label{tab:ketergantungan-mv}
  \begin{tabularx}{\textwidth}{|l|l|l|X|}
    \hline
    \textbf{Varian} & \textbf{Total} & \textbf{View butuh} & \textbf{Cara pengujian} \\
    \hline
    MVC & 3 & Model & Membutuhkan View asli \\
    \hline
    MVP & 2 & tidak ada & Cukup View tiruan \\
    \hline
    MVVM & 3 & ViewModel & Tanpa View, cukup membaca propertinya \\
    \hline
    MVI & 0 & tidak ada & Tanpa View, cukup memanggil fungsi reduksi \\
    \hline
  \end{tabularx}
\end{table}

Angka pada tabel tersebut punya dua bacaan. Pertama, MVI mencatat nol
ketergantungan karena penghubungnya berupa fungsi, bukan kelas yang menyimpan
rujukan ke komponen lain. Kedua, MVVM mencatat total yang sama dengan MVC,
tetapi ketergantungannya menunjuk ke arah yang berbeda. Pada MVC, View mengenal
Model. Pada MVVM, View mengenal ViewModel, sedangkan ViewModel tidak mengenal
View. Jumlah yang sama karena itu tidak berarti sifatnya sama.

\section{Ketika Satu Komponen Menggemuk}

Pemisahan peran hanya bertahan bila setiap komponen menolak pekerjaan yang bukan
miliknya. Bila penolakan itu tidak terjadi, satu komponen menampung hampir semua
kode. Komponen lain tinggal menjadi hiasan. Keadaan tersebut punya nama sendiri,
dan namanya diambil dari komponen yang menggemuk.

\subsection{Fat Controller}

\textit{Fat controller} adalah Controller yang menyimpan aturan bisnis. Nasihat
resminya sederhana. Controller bertugas menerima permintaan, memanggil pihak
yang mengerjakan, lalu menyiapkan balasan~\cite{rails2026controller}. Aturan
bisnisnya bukan miliknya.

Listing~\ref{lst:fat-controller} memperlihatkan bentuk yang keliru. Validasi,
perhitungan, dan penyusunan teks menumpuk di satu tempat.

\begin{lstlisting}[style=PythonStyle, caption={Controller yang menggemuk, sebab
  aturan bisnisnya ikut ditulis di dalamnya}, label=lst:fat-controller]
class Controller:

  def tangani_masukan(self, asal, tujuan, nominal):
    if nominal <= 0:                      # aturan bisnis
      return "Galat: nominal harus positif"
    if nominal > 1_000_000_000:           # aturan bisnis
      return "Galat: nominal melewati batas"
    kurs = self.model.cari_kurs(asal, tujuan)
    hasil = nominal * kurs                # perhitungan bisnis
    return f"Hasil: {hasil:,.2f}"         # penyusunan tampilan
\end{lstlisting}

Akibatnya ada tiga. Pertama, aturan tersebut tidak dapat dipakai ulang. Pintu
masuk lain seperti antarmuka baris perintah atau penjadwal harus menyalin ulang
aturan yang sama. Kedua, pengujiannya menuntut konteks permintaan, padahal yang
diuji hanyalah perkalian dan perbandingan angka. Ketiga, aturan yang sama
tersebar di banyak Controller, sehingga satu perubahan kebijakan menuntut
penyuntingan di banyak tempat.

Perbaikannya adalah memindahkan kedua pemeriksaan dan perkaliannya ke Model,
lalu menyisakan pemanggilan saja di Controller. Bab~\ref{bab:lapisan} sudah
memakai pola tersebut pada \texttt{hitung\_konversi}.

Pada pemrograman aplikasi iOS, gejala yang sama dikenal dengan plesetan
\textit{Massive View Controller}. Singkatan MVC dibaca ulang karena View dan
Controller menyatu, lalu keduanya menampung seluruh kode.

\subsection{Fat Model}

\textit{Fat model} adalah kebalikannya. Nasihat \textit{skinny controller, fat
model} memang benar sebagai arah, tetapi sering dibaca berlebihan. Model lalu
menampung empat hal sekaligus, yaitu aturan bisnis, akses basis data, penyusunan
teks tampilan, dan pengiriman pemberitahuan.

Model seperti itu punya banyak alasan untuk berubah. Perubahan tata letak
kolom basis data menyentuhnya. Perubahan format tanggal menyentuhnya. Perubahan
kebijakan diskon juga menyentuhnya. Kelas dengan ribuan baris biasanya lahir
dari sini.

Perbaikannya bukan mengembalikan kode ke Controller. Perbaikannya adalah membagi
Model menjadi beberapa bagian, misalnya bagian aturan domain, bagian akses
penyimpanan, dan bagian penyusunan teks. Bab~\ref{bab:port} membahas pemisahan
tersebut lebih jauh memakai port dan adapter.

\subsection{Fat View}

\textit{Fat view} adalah View yang menyimpan logika. Gejalanya berupa
percabangan dan perhitungan di dalam templat. Contohnya penghitungan diskon,
penentuan peringkat, atau pemilihan pesan berdasarkan banyak syarat.

Kerugiannya jelas. Logika di dalam templat sulit diuji tanpa menggambar
antarmuka. Logika tersebut juga mudah tersalin ke templat lain, sebab menyalin
lebih cepat daripada memindahkannya ke tempat yang benar. MVP dan MVI menekan
gejala ini secara bawaan, sebab View pada kedua varian tersebut hanya menerima
hasil yang sudah jadi.

Tabel~\ref{tab:gemuk} merangkum ketiga gejala tersebut.

\begin{table}[htbp]
  \centering
  \caption{Gejala komponen yang menggemuk beserta perbaikannya}
  \label{tab:gemuk}
  \begin{tabularx}{\textwidth}{|l|X|X|}
    \hline
    \textbf{Gejala} & \textbf{Tanda yang terlihat} & \textbf{Perbaikan} \\
    \hline
    Fat controller & Validasi dan perhitungan ada di Controller & Pindahkan aturannya ke Model \\
    \hline
    Fat model & Satu kelas mengurus aturan, basis data, dan format & Bagi menjadi beberapa bagian menurut alasan berubahnya \\
    \hline
    Fat view & Percabangan dan perhitungan ada di templat & Susun teksnya di Presenter atau di fungsi reduksi \\
    \hline
  \end{tabularx}
\end{table}

Satu catatan penting menutup bagian ini. Alat \texttt{periksa\_mv.py} pada bab
ini mengukur arah ketergantungan, bukan ukuran komponennya. Sebuah program bisa
saja mencatat angka yang rapi, tetapi seluruh aturannya menumpuk di satu kelas.
Arah panah yang benar karena itu belum menjamin pembagian kerja yang sehat.
Jumlah baris dan jumlah method perlu ikut dilihat.

\section{Kelebihan dan Kekurangan}

Tabel~\ref{tab:untung-rugi-mv} merangkum untung rugi keempat varian tersebut.

\begin{table}[htbp]
  \centering
  \caption{Perbandingan keempat varian Model-View-*}
  \label{tab:untung-rugi-mv}
  \begin{tabularx}{\textwidth}{|l|X|X|}
    \hline
    \textbf{Varian} & \textbf{Kelebihan} & \textbf{Kekurangan} \\
    \hline
    MVC & Paling sederhana, banyak didukung framework & View mengenal Model, sehingga sulit diuji terpisah \\
    \hline
    MVP & View pasif dan mudah diganti tiruan & Presenter mudah membengkak karena memuat seluruh penyusunan teks \\
    \hline
    MVVM & Pembaruan tampilan berjalan otomatis lewat pengamatan & Membutuhkan mekanisme properti yang dapat diamati \\
    \hline
    MVI & Keadaan tetap, riwayat dapat diputar ulang & Setiap tindakan harus dibungkus menjadi Intent, sehingga kode bertambah \\
    \hline
  \end{tabularx}
\end{table}

Pemilihan varian jarang bebas sepenuhnya. Framework antarmuka biasanya sudah
menganjurkan satu varian tertentu, dan melawan anjuran tersebut berbiaya mahal.
Yang tetap dapat dipilih adalah seberapa banyak logika yang dititipkan ke View.
Semakin sedikit logika di dalam View, semakin besar bagian aplikasi yang dapat
diuji tanpa menjalankan antarmuka.

\section{Ringkasan}

Kelompok pola Model-View-* memisahkan tampilan dari aturan bisnis. Seluruh
varian membagi aplikasi menjadi bagian yang menyimpan data, bagian yang
menampilkan, dan bagian yang menghubungkan keduanya. Perbedaan antar varian
terletak pada cara bagian penghubung bekerja, bukan pada tujuannya.

Pembeda yang paling menentukan adalah komponen yang dikenal oleh View. Pada MVC
View mengenal Model, pada MVVM View mengenal ViewModel, sedangkan pada MVP dan
MVI View tidak mengenal apa pun. Semakin sedikit yang dikenal View, semakin
besar bagian aplikasi yang dapat diuji tanpa menjalankan antarmuka.

Pengukuran pada bab ini menunjukkan jumlah ketergantungan saja tidak cukup untuk
menilai sebuah varian. MVC dan MVVM sama-sama mencatat tiga ketergantungan,
tetapi arahnya berbeda, dan perbedaan arah itulah yang menentukan kemudahan
pengujiannya. Angka perlu dibaca bersama arahnya.

MVI mencatat nol ketergantungan karena penghubungnya berupa fungsi murni.
Sifat murni tersebut memberi keuntungan tambahan, yaitu seluruh keadaan yang
pernah digambar dapat disimpan lalu diputar ulang. Harganya, setiap tindakan
pengguna harus dibungkus menjadi Intent, sehingga jumlah kode bertambah.

Pemilihan varian pada praktiknya lebih sering ditentukan oleh framework yang
dipakai daripada oleh pertimbangan bebas. Yang tetap dapat dikendalikan adalah
banyaknya logika yang dititipkan ke View. Keputusan itu diambil dengan melihat
angka ketergantungan beserta arahnya, bukan dengan mengikuti kebiasaan.

Nama pola perlu dibaca dengan hati-hati. Django menyebut polanya MTV, Android
menyebut kelas penghubungnya \textit{state holder}, dan Redux mewarisi
arsitektur Elm. Kelas \texttt{Intent} pada Android justru memakai gagasan yang
sama dengan Intent pada MVI, yaitu niat yang dibungkus menjadi objek data.
Kesamaan gagasan tersebut tetap tidak berarti Android memakai MVI, sebab
tingkat penerapannya berbeda. Sebutan yang beredar luas karena itu perlu
diperiksa terhadap dokumentasi resminya.

Arah ketergantungan yang benar juga belum cukup. Satu komponen tetap dapat
menampung hampir semua kode, dan gejalanya dikenal sebagai \textit{fat
controller}, \textit{fat model}, atau \textit{fat view}. Ukuran komponen karena
itu perlu dilihat bersama arah panahnya.

\section{Latihan}

Ketiga latihan berikut memakai satu basis kode yang sama di
\berkas{sources/chapter03/}, tetapi menyasar tiga masalah yang berbeda.
Latihan~1 memeriksa siapa mengenal siapa. Latihan~2 memeriksa kemudahan
pengujian tanpa antarmuka. Latihan~3 mengukur luas dampak satu permintaan
perubahan. Penilaian didasarkan pada kelengkapan tabel hasil dan ketepatan
jawaban analisisnya. Hasil setiap mahasiswa dapat berbeda, dan yang dilaporkan
adalah pola perbandingannya.

\subsection*{Latihan 1: Memetakan Ketergantungan Antar Komponen}

Latihan ini membuktikan bahwa keempat varian berbeda pada arah
ketergantungannya, bukan hanya pada nama komponennya. Langkah kerjanya sebagai
berikut.

\begin{enumerate}
  \item Jalankan skrip pemeriksaan seperti pada Listing~\ref{lst:jalankan-periksa}.
  \item Catat baris \texttt{Total ketergantungan} dan \texttt{Ketergantungan
    View} untuk keempat varian.
  \item Buka keempat berkas varian, lalu tunjukkan baris konstruktor yang
    menjadi sumber setiap ketergantungan tersebut.
  \item Gambarkan arah panah antar komponen untuk keempat varian.
\end{enumerate}

\begin{lstlisting}[language=bash, caption={Menjalankan pemeriksaan ketergantungan},
  label=lst:jalankan-periksa]
cd sources/chapter03
python periksa_mv.py
\end{lstlisting}

Tabel~\ref{tab:lat1mv-contoh} memperlihatkan hasil satu kali pemeriksaan nyata.

\begin{table}[htbp]
  \centering
  \caption{Contoh hasil Latihan 1 yang sudah terisi}
  \label{tab:lat1mv-contoh}
  \begin{tabularx}{\textwidth}{|l|l|l|X|}
    \hline
    \textbf{Varian} & \textbf{Total} & \textbf{View butuh} & \textbf{Komponen lain} \\
    \hline
    MVC & 3 & 1 & Controller butuh Model dan View \\
    \hline
    MVI & 0 & 0 & Penghubungnya berupa fungsi \\
    \hline
  \end{tabularx}
\end{table}

Isikan hasil pemeriksaan pada Tabel~\ref{tab:lat1mv-hasil}.

\begin{table}[htbp]
  \centering
  \caption{Lembar isian Latihan 1}
  \label{tab:lat1mv-hasil}
  \begin{tabularx}{\textwidth}{|l|l|l|X|}
    \hline
    \textbf{Varian} & \textbf{Total} & \textbf{View butuh} & \textbf{Komponen lain} \\
    \hline
    MVC & \dots & \dots & \dots \\
    \hline
    MVP & \dots & \dots & \dots \\
    \hline
    MVVM & \dots & \dots & \dots \\
    \hline
    MVI & \dots & \dots & \dots \\
    \hline
  \end{tabularx}
\end{table}

Jawab dua pertanyaan berikut dengan menunjuk angka pada tabel. Pertama, MVC dan
MVVM mencatat total yang sama, sehingga jelaskan mengapa keduanya tetap berbeda
sifatnya. Kedua, jelaskan penyebab MVI mencatat nol ketergantungan.

\subsection*{Latihan 2: Menguji Tanpa Antarmuka}

Latihan ini membuktikan hubungan antara ketergantungan View dan kemudahan
pengujian. Langkah kerjanya sebagai berikut.

\begin{enumerate}
  \item Jalankan \texttt{python uji\_tanpa\_antarmuka.py}, lalu catat cara uji
    yang dipakai setiap varian.
  \item Buka \berkas{sources/chapter03/uji_tanpa_antarmuka.py}, lalu bandingkan
    isi keempat fungsi ujinya.
  \item Coba tulis ulang pengujian MVC tanpa membuat objek View, lalu catat apa
    yang menghalangi.
  \item Hitung berapa baris kode tambahan yang dibutuhkan setiap varian untuk
    dapat diuji.
\end{enumerate}

Tabel~\ref{tab:lat2mv-contoh} memperlihatkan hasil satu kali pengujian nyata.

\begin{table}[htbp]
  \centering
  \caption{Contoh hasil Latihan 2 yang sudah terisi}
  \label{tab:lat2mv-contoh}
  \begin{tabularx}{\textwidth}{|l|l|X|}
    \hline
    \textbf{Varian} & \textbf{Lulus} & \textbf{Cara uji} \\
    \hline
    MVC & ya & Membutuhkan View asli \\
    \hline
    MVI & ya & Tanpa View, cukup memanggil fungsi reduksi \\
    \hline
  \end{tabularx}
\end{table}

Isikan hasil pengujian pada Tabel~\ref{tab:lat2mv-hasil}.

\begin{table}[htbp]
  \centering
  \caption{Lembar isian Latihan 2}
  \label{tab:lat2mv-hasil}
  \begin{tabularx}{\textwidth}{|l|l|X|}
    \hline
    \textbf{Varian} & \textbf{Lulus} & \textbf{Cara uji dan baris tambahan} \\
    \hline
    MVC & \dots & \dots \\
    \hline
    MVP & \dots & \dots \\
    \hline
    MVVM & \dots & \dots \\
    \hline
    MVI & \dots & \dots \\
    \hline
  \end{tabularx}
\end{table}

Jawab dua pertanyaan berikut. Pertama, cocokkan kolom cara uji pada latihan ini
dengan kolom ketergantungan View pada Latihan~1, lalu jelaskan hubungan antara
keduanya. Kedua, sebutkan varian yang paling murah diuji beserta alasannya
menurut angka yang diperoleh.

\subsection*{Latihan 3: Mengukur Luas Dampak Perubahan}

Latihan ini membuktikan berapa banyak berkas yang tersentuh oleh satu permintaan
perubahan yang sama. Permintaannya adalah menambahkan baris kedua pada tampilan
yang menyebutkan nilai kurs yang dipakai. Langkah kerjanya sebagai berikut.

\begin{enumerate}
  \item Catat lebih dahulu berkas apa saja yang menurut perkiraan akan berubah
    pada setiap varian.
  \item Terapkan perubahan tersebut pada keempat varian, satu per satu.
  \item Jalankan kembali \texttt{python uji\_tanpa\_antarmuka.py} untuk
    memastikan keempatnya masih berjalan.
  \item Hitung berkas, kelas, dan baris yang berubah pada setiap varian memakai
    perintah pada Listing~\ref{lst:hitung-diff}.
\end{enumerate}

\begin{lstlisting}[language=bash, caption={Menghitung baris yang berubah},
  label=lst:hitung-diff]
cd sources/chapter03
git diff --stat mvc.py mvp.py mvvm.py mvi.py
\end{lstlisting}

Tabel~\ref{tab:lat3mv-contoh} memperlihatkan bentuk pengisian yang diharapkan,
memakai angka perkiraan sebagai contoh.

\begin{table}[htbp]
  \centering
  \caption{Contoh bentuk pengisian Latihan 3}
  \label{tab:lat3mv-contoh}
  \begin{tabularx}{\textwidth}{|l|l|l|X|}
    \hline
    \textbf{Varian} & \textbf{Kelas} & \textbf{Baris} & \textbf{Komponen yang berubah} \\
    \hline
    MVC & 2 & 6 & Model dan View \\
    \hline
    MVI & 2 & 5 & Keadaan dan fungsi reduksi \\
    \hline
  \end{tabularx}
\end{table}

Isikan hasil pengukuran pada Tabel~\ref{tab:lat3mv-hasil}.

\begin{table}[htbp]
  \centering
  \caption{Lembar isian Latihan 3}
  \label{tab:lat3mv-hasil}
  \begin{tabularx}{\textwidth}{|l|l|l|X|}
    \hline
    \textbf{Varian} & \textbf{Kelas} & \textbf{Baris} & \textbf{Komponen yang berubah} \\
    \hline
    MVC & \dots & \dots & \dots \\
    \hline
    MVP & \dots & \dots & \dots \\
    \hline
    MVVM & \dots & \dots & \dots \\
    \hline
    MVI & \dots & \dots & \dots \\
    \hline
    Perkiraan awal yang meleset & \multicolumn{3}{l|}{\dots} \\
    \hline
  \end{tabularx}
\end{table}

Jawab tiga pertanyaan berikut dengan menunjuk angka pada tabel. Pertama, varian
mana yang paling sedikit tersentuh, dan apakah hasilnya sesuai perkiraan awal.
Kedua, jelaskan mengapa varian dengan ketergantungan paling sedikit belum tentu
yang paling sedikit berubah. Ketiga, dengan menggabungkan hasil ketiga latihan,
tentukan varian yang dipilih untuk aplikasi tersebut beserta angka yang menjadi
dasar keputusannya.
